% ASTRO-NOT — THE 2050 PROJECT: production explainer
% Compile: pdflatex project-explained.tex (run twice for TOC/refs)
\documentclass[11pt,a4paper]{article}
\usepackage[margin=22mm]{geometry}
\usepackage[T1]{fontenc}
\usepackage[utf8]{inputenc}
\usepackage{lmodern}
\usepackage{amssymb}
\usepackage{xcolor}
\usepackage{longtable}
\usepackage{enumitem}
\usepackage{listings}
\usepackage{fancyhdr}
\usepackage{hyperref}
\usepackage{titlesec}

\hypersetup{colorlinks=true, linkcolor=linknavy, urlcolor=linkcyan,
  pdftitle={ASTRO-NOT -- The 2050 Project: How This Magazine Was Built},
  pdfauthor={ASTRO-NOT production team}}

\definecolor{navy}{HTML}{101B33}
\definecolor{deepblue}{HTML}{193F68}
\definecolor{icycyan}{HTML}{83DBF3}
\definecolor{paper}{HTML}{F4F6F3}
\definecolor{sun}{HTML}{F3A86C}
\definecolor{linknavy}{HTML}{193F68}
\definecolor{linkcyan}{HTML}{1B6C8F}
\definecolor{codebg}{HTML}{F4F6F3}

\titleformat{\section}{\Large\bfseries\color{navy}}{\thesection}{0.6em}{}
\titleformat{\subsection}{\large\bfseries\color{deepblue}}{\thesubsection}{0.6em}{}
\titleformat{\subsubsection}{\normalsize\bfseries\color{deepblue}}{\thesubsubsection}{0.6em}{}

\lstset{basicstyle=\ttfamily\small, backgroundcolor=\color{codebg},
  frame=single, rulecolor=\color{deepblue}, framerule=0.4pt,
  breaklines=true, columns=fullflexible, showstringspaces=false,
  keywordstyle=\bfseries, commentstyle=\itshape\color{deepblue},
  literate={→}{{$\rightarrow$}}1}

\pagestyle{fancy}
\fancyhf{}
\fancyhead[L]{\small\color{deepblue}ASTRO-NOT $\cdot$ THE 2050 PROJECT}
\fancyhead[R]{\small\color{deepblue}\thepage}
\renewcommand{\headrulewidth}{0.4pt}

\newcommand{\badge}[1]{\textcolor{deepblue}{\textbf{[#1]}}}
\newcommand{\code}[1]{\texttt{#1}}

\begin{document}

\begin{titlepage}
\centering
\vspace*{18mm}
{\color{deepblue}\Large THE 2050 PROJECT}\\[4mm]
{\color{navy}\Huge\bfseries ASTRO-NOT}\\[6mm]
{\color{navy}\LARGE\bfseries How This 48-Page Magazine Was Built}\\[8mm]
{\color{deepblue}\large A production explainer: pipeline, design system,\\ editorial rules, QA and how to rebuild everything}\\[14mm]
{\color{sun}\rule{80mm}{1.2pt}}\\[10mm]
\begin{minipage}{0.8\textwidth}\centering\color{deepblue}
\textbf{Thesis of the magazine it documents:} the Universe will \emph{not}
change drastically by 2050 --- our instruments, discoveries and presence in
space may transform our \emph{understanding} of it. Every claim is labelled,
every simulation marked, every source registered.
\end{minipage}\\[14mm]
{\color{sun}\rule{80mm}{1.2pt}}
\vfill
{\small\color{deepblue}astro-not-magazine/docs/project-explained.tex $\cdot$ rebuilt by \code{build.py} $\cdot$ today's date on the cover of the magazine itself}
\end{titlepage}

\tableofcontents
\newpage

% ==================================================================
\section{What was built, in one page}

\textbf{ASTRO-NOT / THE 2050 PROJECT} is a 48-page, full-colour A4 science
magazine arguing a disciplined editorial thesis: on 24-year timescales the
Universe itself is stable --- its laws, its expansion, its inventory --- while
our \emph{picture} of it, our instruments, and humanity's footprint in space
can change substantially. The magazine exists in two editions, produced
entirely by scripts in this repository:

\begin{itemize}[leftmargin=6mm,itemsep=1pt]
  \item \textbf{Print edition} --- \code{output/astronot-print.pdf}: A4 trim
        $210\times297$~mm with a 3~mm bleed on every edge (page box
        $216\times303$~mm), 15~mm safe margins, three-column flexible grid,
        alternating dark/light spreads, embedded open-licensed fonts.
  \item \textbf{Digital edition} --- \code{output/astronot-digital.pdf}: plain
        A4, tagged PDF (PDF/UA-1 attempt via WeasyPrint), alt text on every
        figure, clickable internal links from page references to page anchors,
        page numbers on every sheet.
\end{itemize}

Supporting both editions:
\begin{itemize}[leftmargin=6mm,itemsep=1pt]
  \item \textbf{Six interactive missions} (browser activities, M01--M06) in
        \code{missions/}, each with a verified answer key and automated tests.
  \item \textbf{Eleven original seeded figures} in \code{figures/out/} ---
        reproducible by a fixed random seed, each simulated one carrying a
        visible \textbf{SIMULATED} banner, axis labels, units and uncertainties.
  \item \textbf{Seven QR codes} generated from one \code{qr/links.yaml} register.
  \item \textbf{Three registers} that make the honesty auditable: references
        (\code{registers/references.yaml}), assets/licences
        (\code{assets/asset-register.yaml}) and links (\code{qr/links.yaml}).
  \item \textbf{A QA gate} (\code{qa/qa\_check.py}) that enforces the
        non-negotiable editorial rules programmatically.
\end{itemize}

\subsection*{The one-command rebuild}
\begin{lstlisting}
.venv/bin/python build.py      # or: make all
# stages: tokens -> figures -> QR -> mission tests
#         -> layout (print+digital PDFs) -> CMYK attempt -> QA gate
\end{lstlisting}
First-time setup (PEP 668 blocks system pip on this machine, so a virtual
environment is used):
\begin{lstlisting}
python3 -m venv .venv
.venv/bin/pip install matplotlib weasyprint pyyaml segno markdown
\end{lstlisting}

% ==================================================================
\section{Repository structure}

\begin{longtable}{p{44mm}p{102mm}}
\textbf{Path} & \textbf{What it is / how it works}\\
\endhead
\code{manuscript/} & 48 Markdown files, one per page. Each begins with YAML
front matter: \code{page}, \code{headline}, \code{section}, \code{type}
(cover / opener / article / infographic / mission / reference), \code{dark}
(bool, drives the dark-or-light spread), \code{confidence} (labels used),
\code{references} (IDs it must cite), and for figure pages: \code{figure},
\code{figure\_alt}, \code{simulated}, \code{obs\_id}.\\
\code{design/tokens.py} & Single source of truth. Defines the palette
(\#101B33 navy, \#193F68 deep blue, \#83DBF3 cyan, \#F4F6F3 paper, \#F3A86C
solar accent), the four confidence labels with their icons and colours, page
geometry (trim, bleed, margins, columns, type sizes) and fonts. Running it
writes \code{tokens.css} and \code{tokens.json} --- nothing downstream may
hardcode a colour.\\
\code{figures/make\_figures.py} & Matplotlib, seeded RNG
(\code{numpy.random.default\_rng(20500117)}) so every run is byte-identical.
Writes SVG. Simulated plots get the red \textbf{SIMULATED} banner
(\code{save(fig, name)} vs \code{save(fig, name, simulated=False)} for
non-simulated charts such as the clock-precision timeline).\\
\code{qr/links.yaml + make\_qr.py} & All seven QR destinations declared once;
\code{segno} writes SVGs into \code{qr/out/}. Local mission QRs point at
relative paths; the portal QR is an explicit \code{.example} placeholder
domain --- deliberately not an invented URL.\\
\code{missions/} & Six self-contained HTML activities + shared
\code{style.css}. Each has a pure-logic \code{logic.js} (no DOM access, so
\code{node} can execute it) and an \code{index.html} that draws on canvas.
\code{missions/tests/test\_missions.py} runs the logic under Node and checks
HTML requirements.\\
\code{layout/} & \code{build\_layout.py}: Markdown $\rightarrow$ HTML
(python-markdown) $\rightarrow$ inline CSS $\rightarrow$ WeasyPrint PDFs.
\code{style.css} + \code{style-print.css}/\code{style-digital.css}.
\code{cmyk\_attempt.py}: Ghostscript colour-conversion attempt.\\
\code{registers/references.yaml} & R01--R22: organisation, date, title,
locator, \emph{locator\_verified} flag and page-use notes. Nothing was
invented; uncertain DOIs are tagged \texttt{[NEEDS VERIFICATION]}.\\
\code{assets/asset-register.yaml} & Every asset with its licence and status.
Photographs are deliberately absent: any NASA/ESA/ESO candidate needs
individual licence review (including commercial print rights) before
download, so original vector art is used meanwhile.\\
\code{qa/} & \code{qa\_check.py} (the rule gate) and \code{link\_checker.py}
(QR targets + internal anchors).\\
\code{build.py, Makefile} & Orchestrator; exits non-zero at the first failing
stage, so a red build means a real rule was violated.\\
\code{output/} & The two PDFs, HTML intermediates, and the CMYK proof attempt.\\
\end{longtable}

% ==================================================================
\section{Toolchain decisions, honestly recorded}

\subsection{Why not Scribus}
The blueprint prescribed Scribus headless (\code{scribus -g -py}). Scribus was
not installed, and headless Scribus scripting is fragile and poorly documented
for automated builds. \textbf{Fallback chosen: HTML/CSS $\rightarrow$
WeasyPrint.} Recorded in \code{BUILD\_LOG.md}. This route gives exact CSS
control over page boxes, bleeds, column grids and dark/light backgrounds, and
runs in $\sim$14~s for both editions.

\subsection{WeasyPrint quirks discovered}
\begin{enumerate}[leftmargin=6mm,itemsep=1pt]
  \item \code{@page \{ size: 216mm 303mm \}} passed via the \code{css=}
        parameter was silently ignored; the same rule inlined in a
        \code{<style>} element works. All CSS is now inlined.
  \item Page numbers were NOT taken from \code{@page} margin boxes
        (\code{counter(page)} lives only there, and margin boxes fight with
        full-bleed backgrounds). Instead every page carries a \emph{static}
        footer ``Page NN / 48'' generated from its front matter --- which also
        makes the 1:1 page mapping auditable by text extraction.
  \item One section per sheet is guaranteed by giving each page a unique CSS
        named page (\code{section\#page-N \{ page: pNN; \}}) --- a forced
        break-after alone was not reliable.
\end{enumerate}

\subsection{Fonts}
System-installed, open-licensed only: \textbf{Noto Sans} (SIL OFL 1.1) for
body and headings, \textbf{DejaVu Sans Mono} for readouts, and
\textbf{DejaVu Sans} as fallback because it carries the confidence icons
($\bullet$, $\blacklozenge$, $\blacktriangle$, $\star$, arrows) that Noto Sans
lacks --- without it, WeasyPrint silently substitutes serif faces for those
glyphs. The OFL licence text ships with the bundle. Serif faces still listed
by \code{pdffonts} come from text inside the matplotlib SVGs, not body copy.

\subsection{CMYK / PDF-X}
Ghostscript \code{pdfwrite} with \code{-sColorConversionStrategy=CMYK}
(\code{layout/cmyk\_attempt.py}) produces \code{astronot-print-cmyk.pdf} as a
\emph{proofing} artefact. Honest limitation: this is colour re-encoding, not a
press-ready X-1a with an embedded OutputIntent ICC --- the print shop must do
final separations. Stated in README and BUILD\_LOG rather than overclaimed.

% ==================================================================
\section{The design system}

\subsection{Tokens defined once}
Colours, the four confidence labels, page geometry and font stacks live only
in \code{design/tokens.py}; the CSS is \emph{generated} from it
(\code{tokens.css}). Change a token, rebuild, and both editions update.

\subsection{The confidence-label system}
Every scientific claim is tagged with exactly one label; icons and colours are
fixed system-wide:

\begin{longtable}{p{10mm}p{44mm}p{92mm}}
\textbf{Icon} & \textbf{Label} & \textbf{Meaning / colour}\\
\endhead
$\bullet$ & Established science & verified repeatedly, independently
reproduced --- deep green \#2E7D4F\\
$\blacklozenge$ & Active research & under study; conclusions may change ---
deep blue \#193F68\\
$\blacktriangle$ & Plausible projection & reasoned extrapolation, not a
promise --- amber-brown \#B0722A\\
$\star$ & Fictional scenario & invented for storytelling, always marked ---
violet \#8A4FA3\\
\end{longtable}

In manuscripts these are written inline as \code{::est::}, \code{::act::},
\code{::pla::}, \code{::fic::}; the builder converts them to styled spans.
Page-level badges are generated from the front-matter \code{confidence} list.
Contrast of every text/background pair is checked automatically against WCAG
AA (worst acceptable pair currently 5.05:1 --- pass).

\subsection{Page types and spreads}
Four page types, chosen per page in front matter:
\emph{cinematic chapter opener} (dark, chapter number, accent rule),
\emph{explanatory article} (three-column body, evidence chart),
\emph{infographic-led spread} (tables/diagrams dominate),
\emph{interactive challenge page} (mission tag + QR box). The
\code{dark: true/false} flag alternates spreads; dark pages set
\code{background: \#101B33} full-bleed with cyan headings and light captions.

\subsection{Type system}
Body 10.5~pt, captions 8.5~pt, footers 7.5~pt --- all high contrast, all from
the tokens file.

% ==================================================================
\section{How a page gets built (the pipeline in detail)}

\begin{enumerate}[leftmargin=6mm]
  \item \textbf{Tokens} --- \code{design/tokens.py} regenerates
        \code{tokens.css}/\code{tokens.json}.
  \item \textbf{Figures} --- \code{make\_figures.py} draws all eleven SVGs
        with the seeded RNG. Simulated ones receive the red banner at render
        time; non-simulated ones (e.g.\ the telescope-sensitivity chart, whose
        post-2026 region is explicitly labelled an extrapolation) do not.
  \item \textbf{QR} --- \code{make\_qr.py} reads \code{links.yaml}, writes
        \code{qr/out/Q-*.svg}.
  \item \textbf{Mission tests} --- \code{missions/tests/test\_missions.py}
        executes each \code{logic.js} under Node: M01's depth formula, M02's
        matched-filter detection with a $<10\%$ false-alarm requirement over
        50 seeded noise streams, M03's budget model, M04's 12-case answer key,
        M05's case resolution, M06's storage shape --- plus HTML checks
        (SIMULATED/FICTIONAL/PRIVACY labels present, keyboard controls,
        \code{aria-live} feedback).
  \item \textbf{Layout} --- \code{build\_layout.py} parses each manuscript
        file's front matter, converts Markdown (with tables), applies inline
        confidence tokens and \texttt{[NEEDS VERIFICATION]} highlighting,
        inserts the figure (with alt text + SIMULATED banner if flagged), the
        QR box, the confidence badges and the static footer; assigns each page
        a unique named page; then WeasyPrint renders both editions.
  \item \textbf{CMYK attempt} --- Ghostscript proofing conversion.
  \item \textbf{QA gate} --- \code{qa\_check.py} (next section).
\end{enumerate}

% ==================================================================
\section{The QA gate: rules turned into code}

\code{qa/qa\_check.py} fails the build if any of these is violated:

\begin{longtable}{p{66mm}p{80mm}}
\textbf{Check} & \textbf{Method}\\
\endhead
48 pages numbered 1--48 & parse every manuscript's front matter\\
Claims map to references & every \code{references:} ID must exist in the
register\\
No guessed facts & \texttt{[NEEDS VERIFICATION]} markers counted (11)\\
AN-2050-017 consistency & must appear on exactly pages 12, 41, 42, 43, 44 ---
and in mission M05's HTML \emph{and} logic\\
SIMULATED labelling & simulated figure types must be flagged in front matter;
``SIMULATED'' text must appear in both PDFs (pdftotext)\\
Fictional labelling & mission/fictional pages counted (13)\\
Logo untouched & placeholder exists; register must flag ``BLOCKER --- brand
approval''\\
Page 45 placeholders & five \texttt{PLACEHOLDER} slots required\\
QR register complete & every \code{links.yaml} entry has an SVG; portal links
must be \code{.example} placeholders\\
PDF structure & pdfinfo: 48 pages; page box $612.283\times858.898$~pt (print)
/ $595.276\times841.89$~pt (digital); pdffonts: all fonts embedded\\
Word count & body text (page furniture excluded) within 9{,}000--12{,}000\\
Contrast & WCAG ratio $\ge 4.5$ for all defined pairs\\
\end{longtable}

\code{qa/link\_checker.py} additionally verifies that every local QR target
exists on disk and that all internal anchors in the digital HTML resolve.

% ==================================================================
\section{The missions: design and verification}

All six are static HTML/CSS/vanilla-JS --- zero network dependencies ---
keyboard-friendly (native \code{<input type=range>}, \code{<select>},
\code{<button>} with focus outlines), with \code{aria-live} result regions and
SIMULATED/FICTIONAL/PRIVACY banners in-app. Pure logic is separated into
\code{logic.js} files that export via \code{module.exports} so Node can run
them headlessly; \code{index.html} handles only drawing and interaction.

\begin{longtable}{p{12mm}p{38mm}p{96mm}}
\textbf{ID} & \textbf{Activity} & \textbf{Mechanics / verified answer key}\\
\endhead
M01 & Exoplanet transit (p16) & seeded light curve, depth fixed at 1.56\%;
slider varies $R_p/R_s$; depth $=(R_p/R_s)^2$; test asserts the reference
$\sqrt{0.0156}\approx0.125$ and that the curve dips correctly.\\
M02 & Signal or noise (p20) & seeded stream, 50\% chance of a chirp; matched
filter against a template; threshold 0.35; test requires signal detected,
noise rejected, and $\le 10\%$ false alarms over 50 pure-noise streams.\\
M03 & Lunar base designer (p31) & 100 fictional units across power, shielding,
water, research; hazards: 14-day night (power $\ge 30$), radiation (shielding
$\ge 25$), water ($\ge 20$), morale (research $\ge 15$); over-budget fails;
reference allocation verified.\\
M04 & Human vs AI (p36) & 12 deterministic detections with a fixed answer
key; the simulated classifier has documented confusions (cosmic-ray hits,
diffraction spikes, one irregular transient); test asserts 12 cases, key
self-consistency, and an imperfect AI.\\
M05 & Cosmic detective (pp41--44) & fictional case AN-2050-017; three
hypotheses; the answer portal resolves to H3 (instrumental) with a
four-point rationale; test pins the observation ID and the answer.\\
M06 & Prediction log (p8) & five prompts with confidence sliders and
falsification fields; data kept in \code{localStorage} only --- the UI states
this; JSON export; validator checks storage shape.\\
\end{longtable}

% ==================================================================
\section{Editorial integrity: how each non-negotiable rule is met}

\begin{longtable}{p{8mm}p{138mm}}
\textbf{\#} & \textbf{How it is implemented}\\
\endhead
1 & No fabricated evidence: simulated plots/log/spectra all carry the red
SIMULATED banner, axis labels, units and error bars; the banner text is
defined once in \code{tokens.py}. Simulated figure \emph{front matter} is
flagged so QA can cross-check.\\
2 & No invented contributors/quotes: page 45 carries five explicit
\texttt{PLACEHOLDER} slots plus the submission policy (consent, withdrawal
right). QA requires the placeholders to exist.\\
3 & Logo: \code{assets/logo/PLACEHOLDER\_LOGO.svg} is used verbatim and never
redrawn; the register marks the official artwork + brand approval as a
BLOCKER.\\
4 & Images/licences: nothing was downloaded. All visuals are original seeded
figures or hand-written SVG. The register states that a credit line is not
permission and that NASA/ESA/ESO candidates need individual review including
commercial print rights.\\
5 & Claims$\to$references: front-matter \code{references} must match
registered IDs; unverifiable facts are written as
\texttt{[NEEDS VERIFICATION]} (11 instances) rather than guessed; no invented
URLs/DOIs (uncertain locators flagged \code{locator\_verified: false}).\\
6 & Fictional 2050 events (M03's model, M05's case, the ``city on Mars''
verdict) are labelled \texttt{$\star$ Fictional scenario} wherever they
appear.\\
7 & AN-2050-017: one fictional ID, pinned by QA to pages 12, 41--44 and the
M05 portal, with a consistent spectrum on p12 and the two-night evidence
figure on p42.\\
8 & Indian space science (p29): verified capabilities (Chandrayaan-3 south-pole
landing, Aditya-L1, MOM, launch families) are cited to ISRO's official pages;
ambitions (Gaganyaan, sample return) are labelled \texttt{$\blacktriangle$}/
\texttt{$\blacklozenge$} with targets-not-guarantees language.\\
\end{longtable}

% ==================================================================
\section{What remains open before artwork lock}

\begin{enumerate}[leftmargin=6mm,itemsep=1pt]
  \item Official ASTRO-NOT logo file and brand approval (placeholder in place).
  \item Cover key image: individual licence review of NASA Image \& Video
        Library / SVS / ESA / ESO candidates --- a credit line is not
        permission.
  \item The 11 \texttt{[NEEDS VERIFICATION]} items (journal DOIs for R06--R08,
        R15; NIST record; NASA FY budget document; ISRO Gaganyaan page path).
  \item Page 45: real, permissioned voices.
  \item Real portal domain to replace \code{www.astro-not.example}.
  \item Print-shop CMYK separations (the Ghostscript output is proofing only).
\end{enumerate}

\section*{Quick command reference}
\begin{lstlisting}
.venv/bin/python build.py            # rebuild everything + QA gate
pdftoppm -png -r 100 output/astronot-print.pdf preview   # page previews
xdg-open output/astronot-digital.pdf # read the digital edition
xdg-open missions/m01-transit/index.html                  # try a mission
.venv/bin/python qa/qa_check.py      # run the rule gate alone
.venv/bin/python missions/tests/test_missions.py
\end{lstlisting}

\end{document}
